%! Least Squares Approximations — Unit 4, Lesson 4.3 — Lesson Plan
\documentclass[10pt]{article}
\usepackage{linalg-article}
\usepackage{linalg-boxes}
\usepackage{graphicx}   % -article does NOT load graphicx; needed for thumbnails/screenshots
\graphicspath{{images/}}
\newcommand{\TallMath}[1]{$\displaystyle #1\rule[-1.4em]{0pt}{3.2em}$}
\newcommand{\vv}{\mathbf{v}}
\newcommand{\ww}{\mathbf{w}}
\newcommand{\xx}{\mathbf{x}}
\newcommand{\bb}{\mathbf{b}}
\newcommand{\pp}{\mathbf{p}}
\newcommand{\ee}{\mathbf{e}}
\newcommand{\zero}{\mathbf{0}}
\newcommand{\T}{\mathsf{T}}
\newcommand{\UnitNumberName}{Unit 4: Orthogonality \quad}
\newcommand{\LessonNumberName}{Lesson 4.3: Least Squares Approximations}

\begin{document}
\begin{center}
    {\Large\bfseries \CourseName: \SchoolYear \\
    {\small\bfseries \UnitNumberName \LessonNumberName}}
\end{center}

\begin{tcolorbox}[colback=blush, colframe=burgundy, boxrule=0.9pt, arc=2mm,
    left=2mm, right=2mm, top=1.5mm, bottom=1.5mm]
    \textbf{Primary Objective:} Students can fit a straight line $y=C+Dt$ to data that does \emph{not}
    lie on a line by recognizing the fit as a \textbf{projection}: when $A\xx=\bb$ has no solution, the
    \textbf{least-squares} answer $\hat{\xx}$ solves the \textbf{normal equations} $A^{\T}A\hat{\xx}=A^{\T}\bb$.
    They can read the fitted values $\pp=A\hat{\xx}$ and residuals $\ee=\bb-\pp$, explain that least squares
    minimizes the total squared error $\|\ee\|^2$, and justify \emph{why} the best line is the one whose
    residual is perpendicular to the data columns --- the error lands in $N(A^{\T})$ (Lessons~4.1--4.2).
\end{tcolorbox}

\renewcommand{\arraystretch}{1.4}
\begin{skillbox}[Priority Ideas \& Skills]{goldbox}
\begin{tabularx}{\linewidth}{@{}X|X@{}}
\textbf{Priority Ideas \& Skills} & \textbf{Key Understandings (the ``why'')} \\[2pt]
\hline
\vspace{-6pt}
\begin{itemize}[leftmargin=1.1em, nosep, after=\vspace{2pt}]
  \item Set up $A$, $\xx=(C,D)$, $\bb$ for fitting a line $y=C+Dt$ to data.
  \item Recognize that $A\xx=\bb$ has \textbf{no solution} when the points miss a line.
  \item Solve the \textbf{normal equations} $A^{\T}A\hat{\xx}=A^{\T}\bb$ for the best-fit line.
  \item Compute fitted values $\pp=A\hat{\xx}$ and residuals $\ee=\bb-\pp$; interpret and predict.
\end{itemize}
&
\vspace{-6pt}
\begin{itemize}[leftmargin=1.1em, nosep, after=\vspace{2pt}]
  \item Least squares \emph{is} the Lesson~4.2 projection, applied to data $\bb$.
  \item ``Best'' $=$ smallest total \emph{squared} miss $\|\ee\|^2$ (squares can't cancel).
  \item The residual is perpendicular to the columns: $A^{\T}\ee=\zero$, so $\ee\in N(A^{\T})$.
  \item The column of $1$s is the intercept $C$; it forces $\sum e_i=0$.
\end{itemize}
\end{tabularx}
\end{skillbox}

\begin{skillbox}[{Vocabulary, Concepts, \& Theorems}]{greenbox}
\renewcommand{\arraystretch}{1.3}
\begin{tabularx}{\linewidth}{@{}l X@{}}
\textbf{Overdetermined system} & $A\xx=\bb$ with more equations (data points) than unknowns; usually \emph{no} exact solution. \\
\textbf{Least-squares solution $\hat{\xx}$} & the $(C,D)$ making $A\hat{\xx}$ closest to $\bb$; it solves $A^{\T}A\hat{\xx}=A^{\T}\bb$. \\
\textbf{Normal equations} & $A^{\T}A\hat{\xx}=A^{\T}\bb$ --- the same equations as projecting $\bb$ onto $C(A)$. \\
\textbf{Best-fit line} & $y=C+Dt$ built from $\hat{\xx}$; nearest to the data, need not pass through any point. \\
\textbf{Residual $\ee$} & $\ee=\bb-\pp$, the vertical misses; $A^{\T}\ee=\zero$, and least squares minimizes $\|\ee\|^2$. \\
\end{tabularx}
\end{skillbox}

\begin{fixedskillbox}[Activate Prior Knowledge \& Spiral Review]{blush}
    The Warm-Up rehearses the three skills this lesson builds on:
    \textbf{(1)} reading a value off a line $y=C+Dt$ and computing a \emph{miss} (residual);
    \textbf{(2)} building $A^{\T}A$ and $A^{\T}\bb$ from column dot products (Lesson~4.2); and
    \textbf{(3)} \emph{solving a $2\times2$ system} (the shape of the normal equations). Debrief item~3
    aloud: stress that fitting a line ends in an ordinary linear system --- all the machinery is from
    Units~1--2 plus last lesson's projection.
\end{fixedskillbox}

\begin{skillbox}[Hook]{blush}
    You stretch a spring with $0,1,2,3$ kg and measure its length four times, but the ruler is a little off
    each time --- the four points do \emph{not} land on one line. Ask: \textbf{``Can one straight line pass
    through all four points? If not, what should the \emph{best} line even mean?''} Students see that four
    points rarely lie on a line ($A\xx=\bb$ has no solution), so ``best'' must mean \emph{closest}. Name the
    payoff: the best line's heights $\pp$ are the \textbf{projection} of the data $\bb$ onto the plane of
    possible lines, and the leftover misses $\ee=\bb-\pp$ are the \textbf{residuals}. Then pose the lesson
    question: \textbf{``If best means closest, and closest means error perpendicular, can we reuse last
    lesson's normal equations to fit a line?''}
\end{skillbox}

\begin{skillbox}[Lesson]{blush}
\begin{multicols}{2}
\textbf{1. Too many equations.} Fitting $y=C+Dt$ to four points asks $C+Dt_i=y_i$ four times --- two
unknowns, four equations. In matrix form $A\xx=\bb$ with $A$ a column of $1$s next to a column of the
$t$-values. \textbf{Ask:} why is there usually no solution? Because $\bb$ (the data) rarely lands in the
plane $C(A)$.

\textbf{2. Best $=$ closest $=$ projection.} We cannot solve $A\xx=\bb$, so we find the point $\pp=A\hat{\xx}$
\emph{in} $C(A)$ closest to $\bb$ --- the projection from Lesson~4.2. The best $\hat{\xx}$ is the
\textbf{least-squares solution}.

\textbf{3. The normal equations, again.} Closest means the residual $\ee=\bb-A\hat{\xx}$ is perpendicular
to every column: $A^{\T}\ee=\zero$, i.e. $A^{\T}A\hat{\xx}=A^{\T}\bb$. \emph{Same equation as last lesson,
now fitting a line.} Solve the $2\times2$ system for $(C,D)$.

\columnbreak

\textbf{4. Worked numbers.} Spring data $t=(0,1,2,3)$, $\bb=(3,3,5,9)$: $A^{\T}A=\begin{bsmallmatrix} 4&6\\6&14 \end{bsmallmatrix}$,
$A^{\T}\bb=(20,40)$, so $\hat{\xx}=(2,2)$ and the line is $y=2+2t$. Then $\pp=(2,4,6,8)$ and
$\ee=(1,-1,-1,1)$. \textbf{Ask:} what are the entries of $A^{\T}A$? Sums of $1$, of $t$, of $t^2$ --- built
mechanically from the columns.

\textbf{5. Why squared, and where the error goes.} Least squares minimizes $\|\ee\|^2=\sum e_i^2$: squares
cannot cancel and punish big misses. \textbf{Punchline:} $A^{\T}\ee=\zero$ puts $\ee$ in $N(A^{\T})$, the
orthogonal complement of $C(A)$ (Lesson~4.1) --- and because the first column is all $1$s, $\sum e_i=0$.
The data splits into a part \emph{on} the line ($\pp$) plus a perpendicular \emph{miss} ($\ee$).
\end{multicols}
\end{skillbox}

\begin{skillbox}[Active Monitoring]{redbox}
Circulate during the notes and Tier work. Watch for and correct:
\begin{itemize}[leftmargin=1.2em, nosep]
  \item dropping the column of $1$s --- that column \emph{is} the intercept $C$; without it there is no $C$;
  \item trying to solve $A\xx=\bb$ directly --- it has \emph{no} solution, which is the whole reason we project;
  \item adding raw residuals (they cancel to $0$) instead of \emph{squaring} them when talking about ``smallest error.''
\end{itemize}
Cold-call prompts: ``Why does $A\xx=\bb$ have no solution here?'' ``What does a residual measure?'' ``Which
subspace does the error live in, and why does that make $\pp$ closest?'' ``What is $A^{\T}A$ counting?''
\end{skillbox}

\begin{skillbox}[Group Work \& Differentiation]{redbox}
\textbf{Best Fit by Least Squares} activity (tiered; mirrors the notes).
\begin{multicols}{3}
\textbf{Tier R --- Remediate.} Fit a \emph{horizontal} line $y=C$ to three readings (one column of $1$s):
discover that the best constant is the average.

\columnbreak
\textbf{Tier A --- Approaching.} Fit a line $y=C+Dt$ to three points: build $A^{\T}A$, $A^{\T}\bb$, solve,
and check $\ee\perp$ both columns.

\columnbreak
\textbf{Tier E --- Extension.} Fit a line to four points, \emph{predict} a future value, and interpret ---
which subspace holds $\ee$, and how far should you trust an extrapolation?
\end{multicols}
\end{skillbox}

\begin{skillbox}[Individual Work \& Assessment]{redbox}
\textbf{Exit Ticket} (independent, no notes): fit a line to three points via the normal equations and state
$y=C+Dt$; a \textbf{conceptual/justification check} --- say what the residual measures and what least
squares minimizes; and explain \emph{why} we solve the normal equations instead of $A\xx=\bb$. Sort
responses into ``sets up $A^{\T}A\hat{\xx}=A^{\T}\bb$ / tries to solve $A\xx=\bb$ / drops the $1$s column''
to decide whether 4.4 opens with a quick re-anchor on projection.
\end{skillbox}

\begin{skillbox}[Reinforcement \& Extension]{goldbox}
\textbf{Homework:} fit a best constant (mean), a line through three points, and a line through four points
with a prediction; explain why a line usually cannot pass through every point. \textbf{Extension:} verify
$A^{\T}\ee=\zero$ (so $\ee\in N(A^{\T})$, Lesson~4.1) and show the column of $1$s forces $\sum e_i=0$.
\textbf{Preview:} Lesson~4.4, \emph{Orthogonal Matrices \& Gram--Schmidt} --- when the columns of $A$ are
perpendicular unit vectors, $A^{\T}A=I$ and the normal equations become trivial.
\end{skillbox}
\end{document}
